\documentclass[10pt, a4paper]{article}

% Parametri che modificano il file main.tex
% Le uniche parti da cambiare su main.tex sono:
% - vari \vspace tra sezioni
% - tabella azioni da intraprendere
% - sezione altro

\def\data{2024-03-12}
\def\oraInizio{14:00}
\def\oraFine{15:30}
\def\luogo{Piattaforma Discord}

\def\tipoVerb{Interno} % Interno - Esterno

\def\nomeResp{Orlandi G.} % Cognome N.
\def\nomeVer{Bresolin G.} % Cognome N.
\def\nomeSegr{Michelon R.} % Cognome N.

\def\nomeAzienda{Azzurro Digitale}
\def\firmaAzienda{azzurrodigitale.png}
\def\firmaResp{giacomo.png} % nome Responsabile

\def\listaPartInt{
Bresolin G.,
Campese M.,
Ciriolo I.,
Dugo A.,
Feltrin E.,
Michelon R.,
Orlandi G.
}

\def\listaPartEst{
Azzurro B.,
Digitale C.,
}

% Se nessuna revisione: \def\listaRevisioneAzioni {x}
\def\listaRevisioneAzioni {x}

\def\listaOrdineGiorno {
{Aggiornamento da parte di tutti i membri del gruppo riguardo al lavoro di codifica iniziato dopo la riunione precedente;},
{Preparazione per la riunione con l'azienda pianificata in data 2024-03-13;},
{Assegnazione dei task relativi alla parte di testing.}
}

\def\listaDiscussioneInterna {
{Il team ha eseguito un aggiornamento generale relativo alla codifica del backend in corso da una settimana. Considerati i buoni risultati il gruppo ha discusso della possibilità di presentare il lavoro svolto durante la riunione con AzzurroDigitale pianificata per il 2024-03-13;},
{Il team ha riportato la discussione tenuta in settimana per la divisione del lavoro da svolgere riguardante la parte di test di unità delle componenti sviluppate durante la prima parte di codifica del progetto. La suddivisione delle task viene quindi riportata in 'Azioni da intraprendere'.
}
}



% Se nessuna decisione: \def\listaDecisioni {x}
\def\listaDecisioni{
{Il team presenterà all'azienda il lavoro di codifica delle componenti di backend svolto fino ad ora;},
{Il team continuerà la produzione della parte di test di unità delle componenti di backend secondo la suddivisione dei compiti riportata in 'Azioni da intraprendere'.}
}

\usepackage{style}
\usepackage{headerfooter}
\usepackage{caption}
\usepackage{hyperref}
\usepackage{svg}
\usepackage{comment}
\usepackage{setspace}
\definecolor{primarycolor}{RGB}{248, 182, 143}

\title{\titolo}
\author{SWEetCode}

\begin{document}

% PRIMA PAGINA
\begin{titlepage}
    \thispagestyle{empty}
    \begin{tikzpicture}[remember picture, overlay]
        % TRIANGOLI
        \draw[fill=secondarycolor, secondarycolor] (current page.north west) -- (current page.south west) -- (8.8, -28);
        \draw[fill=primarycolor, primarycolor] (-3, 5) -- (4, -13.6) -- (11, 5);

        % LOGO
        \node [xshift=-5cm, yshift=25cm] (logo) at (current page.south east) {\includegraphics[width=6.5cm]{img/logo.png}};

        % SWEETCODE - DATE
        \node [anchor=north east, align=right, xshift=-1.2cm, yshift=20.5cm, text=black] (sweetcode) at (current page.south east) {\fontsize{32pt}{36pt}\selectfont SWEetCode};
        \draw[line width=4pt, lightcol] ([xshift=-3cm, yshift=-0.37cm]sweetcode.south west) -- ([yshift=-0.37cm]sweetcode.south east);
        \node [anchor=north east, align=right, xshift=-1.2cm, yshift=18.7cm, text=black] (date) at (current page.south east){\fontsize{24pt}{24pt} \selectfont Verbale \tipoVerb};

        % NOME FILE
        \node [anchor=north east, text width=15cm, align=right, xshift=-1.2cm, yshift=17cm, text=black] (titolo) at (current page.south east){\fontsize{48pt}{48pt}\textbf{\data}};

        % BOX DATI PARTECIPANTI
        \ifthenelse{\equal{\tipoVerb}{Esterno}}{
            \node[anchor=north east, xshift=-1.2cm, yshift=14.5cm, minimum width=8cm] (box) at (current page.south east){};
        }{
            \node[anchor=north east, xshift=-1.2cm, yshift=12.5cm, minimum width=8cm] (box) at (current page.south east){};
        }

        % RESPONSABILE
        \node[anchor=north west, align=left] (dati2) at (box.north west) {\fontsize{15pt}{15pt}\selectfont \textbf{Responsabile}};
        \draw[line width=4pt, lightcol] (dati2.south west) -- ([xshift=8cm]dati2.south west);
        \node[anchor=north west, align=left] (dati21) at (dati2.south west){\fontsize{13pt}{13pt}\selectfont \nomeResp};

        % VERIFICATORE
        \node[anchor=north west, yshift=-1cm, align=left] (dati3) at (dati21.north west) {\fontsize{15pt}{15pt}\selectfont \textbf{Verificatore}};
        \draw[line width=4pt, lightcol] (dati3.south west) -- ([xshift=8cm]dati3.south west);
        \node[anchor=north west, align=left] (dati31) at (dati3.south west){\fontsize{13pt}{13pt}\selectfont \nomeVer};

        % SEGRETARIO DI RIUNIONE
        \node[anchor=north west, yshift=-1cm, align=left] (dati4) at (dati31.north west) {\fontsize{15pt}{15pt}\selectfont \textbf{Segretario di Riunione}};
        \draw[line width=4pt, lightcol] (dati4.south west) -- ([xshift=8cm]dati4.south west);
        \node[anchor=north west, align=left] (dati41) at (dati4.south west){\fontsize{13pt}{13pt}\selectfont \nomeSegr};
        
        % UNIPD - SWE
        \node [xshift=4.4cm, yshift=2.3cm, draw, secondarycolor, text=white] (uni) at (current page.south west) {\fontsize{20pt}{20pt} \selectfont Università di Padova};
        \node [xshift=0.65cm, yshift=0.7cm, draw, secondarycolor, text=white, below=of uni] (corso) {\fontsize{20pt}{20pt}\selectfont Ingegneria del Software};

        % FIRMA AZIENDA
        %\ifthenelse{\equal{\tipoVerb}{Esterno}}{
           % \draw[line width=4pt, lightcol] ([xshift=-1.2cm, yshift=5.2cm]current page.south east) -- ([xshift=-8cm, yshift=5.2cm]current page.south east);
          %  \node [xshift=-4.8cm, yshift=6.3cm] (logo) at (current page.south east) {\includegraphics[width=6cm]{img/firme/\firmaAzienda}};
          %  \node[anchor=north west, xshift=12.9cm, yshift=4.95cm, align=left] at (current page.south west)
           % {\fontsize{13pt}{13pt}\selectfont L'azienda: \nomeAzienda};
       % }

       %  FIRMA
        \draw[line width=4pt, lightcol] ([xshift=-1.2cm, yshift=1.8cm]current page.south east) -- ([xshift=-8cm, yshift=1.8cm]current page.south east);
        \node [xshift=-4.8cm, yshift=2.45cm] (logo) at (current page.south east) {\includegraphics[width=6cm]{img/firme/\firmaResp}};
        \node[anchor=north west, xshift=12.9cm, yshift=1.45cm, align=left] at (current page.south west)
        {\fontsize{13pt}{13pt}\selectfont Il Responsabile: \nomeResp};
        
    \end{tikzpicture}
\end{titlepage}

% REGISTRO DELLE VERSIONI
%Registro in ordine dalla più recente alla meno recente!

{\renewcommand{\arraystretch}{1.5}
\section*{Registro delle versioni}

\begin{xltabular}{\textwidth}{c|c|c|c|X}
\label{tab:long}

\textbf{Versione} & \textbf{Data} & \quantities{\textbf{Responsabile di}\\\textbf{stesura}}& \textbf{Revisore} & \quantities{\textbf{Dettaglio e}\\\textbf{motivazioni}} \\
\endfirsthead

\textbf{Versione} & \textbf{Data} & \quantities{\textbf{Responsabile di}\\\textbf{stesura}}& \textbf{Revisore} & \quantities{\textbf{Dettaglio e}\\\textbf{motivazioni}} \\
\endhead

\multicolumn{5}{r}{{Continua nella pagina successiva}} \\
\endfoot

\endlastfoot

\hline
v2.23.0(1) & $2024-03-20$ & \quantities{Ciriolo I.} & Feltrin E. & Gestione e visualizzazione delle chat.\\
\hline
v2.21.0(1) & $2024-03-19$ & \quantities{Ciriolo I.} & Michelon R. & Prima stesura.\\
\hline
    
\end{xltabular}}
%\newpage
\newpage
%todo modificare questa section per avere sempre un titolo "grande" ma senza numero
\section*{Candidatura per Product Baseline}
Il gruppo SWEetCode è lieto di presentare la propria candidatura per la revisione di avanzamento Product Baseline(PB), nell'impegno di consegnare il progetto da Voi commissionato, denominato \textbf{Knowledge Management AI} (C1), proposto dall'azienda \textbf{AzzurroDigitale}.\\

Per una panoramica delle attività e una consultazione semplificata dei documenti riguardanti l'PB, il team di SWEetCode estende un invito cordiale a visitare il sito:
\begin{center}
\href{https://sweetcode-team.github.io}{\textcolor{blue}{sweetcode-team.github.io}}
\end{center}

\vspace{1em}
In tale spazio si potranno trovare:
\begin{itemize}
    \item (Analisi dei requisiti v3.0.0(0));
    \item (Glossario v3.0.0(0));
    \item (Norme di progetto v3.0.0(0));
    \item (Piano di progetto v3.0.0(0));
    \item (Piano di qualifica v3.0.0(0));
    \item (Manuale utente v3.0.0(0));
    \item I verbali delle riunioni interne tenute i giorni:
        \begin{itemize}
            \item 2024-02-19;
            \item 2024-02-26;
            \item 2024-03-07;
            \item 2024-03-12;
            \item 2024-03-19;
            \item 2024-03-25;
            \item 2024-04-01;
            \item 2024-04-12.
        \end{itemize}
   \item I verbali delle riunioni esterne tenute i giorni:
        \begin{itemize}
            \item 2023-02-28: AzzurroDigitale;
            \item 2023-03-13: AzzurroDigitale;
            \item 2023-03-26: AzzurroDigitale.
        \end{itemize}    
\end{itemize}

Per consultare il repository del progetto, il team SWEetCode estende un invito cordiale a visitare il link:\\
\begin{center}
\href{https://github.com/sweetcode-team/Knowledge_Management_AI}{\textcolor{blue}{github.com/sweetcode-team/Knowledge\_Management\_AI}}
\end{center}

\newpage
\subsection*{Aggiornamento Impegni}
Il team conclude con la Product Baseline lo sviluppo del progetto confermando la data di consegna prevista per il 2024-04-12.
Il team ha rispettato il preventivo iniziale di 12715,00€, arrivando ad un consumo totale di 12450,00€.
Le tabelle orarie sono riportate nel documento (Piano di progetto v3.0.0(0), \S Stima dei costi di realizzazione).

{\renewcommand{\arraystretch}{1.5}
\begin{table}[H]
\centering
\begin{tabularx}{0.54\textwidth}{c|c|c}
\rowcolor{primarycolor}
    \textbf{Revisione} & \textbf{Data}& \textbf{Costi}\\
\hline
    \textbf{Product Baseline} & 2024-04-12 & 12450,00€\\
\end{tabularx}
\end{table}

\vspace{1em}
Il team è composto dai seguenti membri:
\\
\\
% TABELLA (tenere anche arraystretch che aumenta il padding verticale)
{\renewcommand{\arraystretch}{1.5}
\begin{tabularx}{\textwidth}{X|X}
\textbf{Cognome e Nome} & \textbf{Matricola n°} \\
\hline
Bresolin Gianluca & 2034316 \\
\hline
Campese Martina & 1122439 \\
\hline
Ciriolo Irene & 2043682 \\
\hline
Dugo Alberto & 2042382 \\
\hline
Feltrin Emanuele & 2034314 \\
\hline
Michelon Riccardo & 2042341 \\
\hline
Orlandi Giacomo & 1122253 \\
\end{tabularx}}\\

Cordiali saluti,\\

Il gruppo SWEetCode

\end{document}